\section{Introduccion}
Elaboracion gráficas de dos robots diferentes al robot 1\textbf{}.
\subsection{\textbf{Robot 4}}

\begin{figure}[h]
	\centering
	\includegraphics[width=0.9\linewidth]{robot4 cinematica articular.png}
	\caption{Robot 4, cinematica articular}
	\label{fig:diagramaRobot4}
\end{figure}
\begin{figure}[h]
	\centering
	\subfloat[Graficas Robot 4]{%
		\includegraphics[width=1.02\textwidth]{robot4 graficas.png}%
		\label{fig:Robot4graficas}
	}
\subsection{\textbf{Error objetivo Robot 4}}
Solución alcanzada en muestra 1, iteraciones 5, error = 4.192727e-07

Solución alcanzada en muestra 1, iteraciones 6, error = 1.529095e-11\textbf{}.
\end{figure}
\subsection{\textbf{Robot 2}}

\begin{figure}[h]
	\centering
	\includegraphics[width=0.9\linewidth]{robot2 cinematica articular.png}
	\caption{Robot 2, cinematica articular}
	\label{fig:diagramaRobot2}
\end{figure}
\begin{figure}[h]
	\centering
	\subfloat[Graficas Robot 2]{%
		\includegraphics[width=1.02\textwidth]{robot2 graficas.png}%
		\label{fig:Robot2graficas}
	}
	\subsection{\textbf{Error objetivo Robot 2}}
Solución alcanzada en muestra 1, iteraciones 5, error = 1.500378e-13

Solución alcanzada en muestra 1, iteraciones 6, error = 5.320608e-13\textbf{}.
\end{figure}
\pagebreak